\chapter{Einleitung}
\label{kap:einleitung}

\begin{definition}
  Ein Vektor des Raums lautet in der Standardbasis
  $\vec{v} = a\,\ihat + b\,\jhat + c\,\khat$.
\end{definition}

\begin{satz}[Euklidische Norm]
  \label{satz:norm}
  Für jedes $\vec{v} = a\,\ihat + b\,\jhat + c\,\khat$ gilt
  \[
    \lVert \vec{v} \rVert = \sqrt{a^2 + b^2 + c^2}.
  \]
\end{satz}

\begin{proof}
  Zweimal den Satz des Pythagoras anwenden.
\end{proof}

\begin{beispiel}
  Für $\vec{v} = 3\,\ihat + 4\,\jhat$ ist $\lVert \vec{v} \rVert = 5$
  (Satz~\ref{satz:norm}).
\end{beispiel}

\begin{bemerkung}
  Legen Sie weitere Kapitel im Ordner \texttt{kapitel/} an und binden Sie sie mit \verb|\input| ein.
\end{bemerkung}
